% !TEX root = ../Document.tex
\Chapter{RÉSULTATS ET DISCUSSION}\label{sec:Theme3}
Présentation avec CBF et sans CBF

Présentation de la détection avec SAM 3

Listes des figures pour les résultats

\begin{itemize}
    \item Trajectoire de la pointe de la fourchette dans le nuage persistent map;
    \begin{enumerate}
        \item Nuage de points de la scène;
        \item Trajectoire sans CBF
        \item Trajectoire avec CBF
        \item Cible.
    \end{enumerate}
    \item Évolution de $h(q)$ en fonction du temps;
    \begin{enumerate}
        \item $h_{\text{corr}}$;
        \item ligne zéro.
        \item zones où le CBF est actif (en rouge par exemple).
    \end{enumerate}
    \item Distance obstalce-robot en fonction du temps
    \item Décomposition des vitesses en fonction du temps
    \begin{enumerate}
        \item $\lvert dq_{ff} \rvert$;
        \item $\lvert dq_{feedback} \rvert$;
        \item $\lvert dq_{base} \rvert$
        \item $\lvert dq_{cbf_{delta}} \rvert$
        \item $\lvert dq_{escape} \rvert$
        \item $\lvert dq_{safe} \rvert$.
        \item Objectif: Expliquer le comportement du système et montré quand le FLow Matching dirige le robot seul, quand le CBF intervient, quand l'échappement est nécessaire, si la vitesse corrigée devient trop aggressive, etc.
    \end{enumerate}
    \item Effet du filtre passe bas
    \begin{enumerate}
        \item Vitesse avant CBF brute
        \item Vitesse avant CBF filtrée
    \end{enumerate}
    \item Temps de calcul du pipeline. Faire un boxplot (ou un tableau) sur:
    \begin{enumerate}
        \item Temps de détection avec SAM3
        \item Temps de suivi avec SAM2
        \item Temps de mise à jour persistent map
        \item Temps de génération de trajectoire avec Flow Matching
        \item Temps de calcul de la cinématique inverse
        \item Temps de calcul des points critiques
        \item Temps de calcul de la correction de trajectoire
        \item Temps total du pipeline
    \end{enumerate}
    \item Résistance faces aux occlusions et pertes de détection.
    \begin{enumerate}
        \item Détection dispo / indispo dans le temps
        \item Images présentants pertes de détection puis redétection
    \end{enumerate}
    \item NE PAS OUBLIER LA DÉTECTION INITALE DES FORK TEETH
    \item Cas d'échec (comme la figure de Félix).
\end{itemize}

% =================================================================
\section{Protocol expérimental}
% =================================================================
\subsection{Description des scénarios testés}
\subsubsection{Pick and place}
\subsubsection{Robot assisted feeding}
\subsection{Métriques utilisées}
À compléter.


% =================================================================
\section{validation du modèle géométrique de sécurité}
% =================================================================
\subsection{Précision du SDF du robot}
\subsection{Temps d'évaluation du SDF et du gradient}
\subsection{Visualisation des points critiques sélectionnés}
À compléter.

% =================================================================
\section{Qualité de la génération nominale}
% =================================================================
\subsection{Trajectoire générée sans obstacle critique}
\subsection{Cohérence de la trajectoire avec la cible}
À compléter.

% =================================================================
\section{Sécurité et évitement d'obstacle}
% =================================================================
\subsection{Comparaison des trajectoires avec et sans SDF-CBF}
\subsection{Évolution de $h(q)$ et de la distance obstalce-robot}
\subsection{Déviation entre trajectoire nominale et trajectoire sécurisée}
À compléter.

% =================================================================
\section{Analyse des vitesses et de la fluidité}
% =================================================================
\subsection{Décomposition des vitesses articulaires \texorpdfstring{$dq_{nom}$, $dq_{escape}$, $dq_{safe}$}{dq_nom, dq_escape, dq_safe}}
\subsection{Effet de la vitesse d'échappement tangentielle \texorpdfstring{$dq_{escape}$}{dq_escape}}
\subsection{Effet du filtrage sur la commande sécurisée pré SDF-CBF \texorpdfstring{$dq_{safe}$}{dq_safe}}
À compléter.

% =================================================================
\section{Performance et temps de calcul}
% =================================================================
\subsection{Temps de génération de trajectoire}
\subsection{Temps de sélection des points critiques}
\subsection{Temps de correction SDF-CBF}
\subsection{Délais jusqu'au contrôleur}
À compléter.

% =================================================================
\section{Performances globales et limites}
% =================================================================
\subsection{Taux de succès des scénarios}
\subsubsection{Pick and place}
\subsubsection{Robot assisted feeding}
\subsection{Cas d'échec}
À compléter.



% =================================================================
\section{Discussion}
% =================================================================
